\section{Approach}

This chapter brings Chapters 3 and 4 together in two parts.
The first part is an analysis: it takes the six steps of the workflow in Section 3.4 and asks, for each, which parts suit ordinary code, which suit an LLM agent and which should stay with a person.
The second part is a design: it describes a system of agents, tools and records that would do this work alongside an analyst, and shows how that system would be used.

Section 5.2 presents the analysis, in one table and one figure rather than step by step, since Chapter 3 has already described how each step is done today.
Section 5.3 turns its results into the principles the design follows.
Section 5.4 gives an overview of the system, and Section 5.5 follows the workflow through it.
Section 5.6 describes how the analyst and the agents work together: an agent suggests a piece of work with its reasons, the analyst accepts, edits or replies, and the system keeps what it learns from each reply.
Section 5.7 describes the calibration agent, the one part of the system that needs to choose its own next step.
Section 5.8 then follows one study through the whole system: a scenario in which the Strait of Hormuz is closed, of the kind the model's developers analysed in 2026.
Section 5.9 summarises the agents and how each could be tested.

The analysis rests on the documentation, the code, the 2023 dataset and the working session with Lukas Barner described in Section 3.3.
The design is an argument built on that analysis and on the literature in Chapters 2 and 4; it has not been built or tested, which is the task of the master's thesis.
